% ===================================================================
%  اللوحة رقم ١ — المدرسة. --- الثانوية. --- الصف.
%
%  Traduction, faite en 2026, en arabe levantin d'aujourd'hui, sur
%  l'ido de texto/io. CE FICHIER PORTE LES REGLES DE LA COLONNE : les
%  quinze autres y renvoient.
%
% -------------------------------------------------------------------
%  1. UNE TRADUCTION N'EST PAS UNE TRANSCRIPTION
% -------------------------------------------------------------------
%  Les fichiers de texto/io et texto/fr reproduisent un fac-simile :
%  ils suivent ses lignes (\nl), ses coupures de mot (\cc) et ses
%  feuillets (\begin{VUpage}). CETTE COLONNE NE SUIT RIEN DE TOUT
%  CELA. Elle porte les MEMES cles %%K, le meme appareil macro pour
%  macro, les memes \textsuperscript{(n)} dans le meme ordre — mais
%  pas une seule macro de mise en page du fac-simile. Le gras ne
%  couvre que le TERME : la ponctuation et les enclitiques restent
%  dehors.
%
% -------------------------------------------------------------------
%  2. CE QUE CETTE COLONNE A A DEFENDRE, ET CONTRE QUI
% -------------------------------------------------------------------
%  TOUTES LES COLONNES PRECEDENTES SE DEFENDAIENT CONTRE UNE LANGUE
%  ETRANGERE. Le marathi contre le hindi, le persan contre l'ourdou
%  et l'arabe, le gujarati contre le hindi : deux langues distinctes,
%  un alphabet commun. Celle-ci est la premiere dont LES DEUX
%  VOISINES SONT LA MEME LANGUE QU'ELLE. L'arabe standard (« ar ») en
%  est la forme ecrite, celle qu'on apprend a l'ecole et dans
%  laquelle toute main levantine remonte d'elle-meme des qu'elle
%  ecrit ; l'arabe egyptien (« arz ») en est une autre forme parlee,
%  celle que tout le monde entend au cinema. Les deux sont deja dans
%  le releve, a une colonne d'ici.
%
%  LE DANGER N'EST DONC PAS QU'UN MOT ETRANGER SE GLISSE : C'EST QUE
%  LA MAIN CHANGE DE REGISTRE SANS QUE PERSONNE LE REMARQUE. On
%  n'ecrit pas « ليس » par erreur : on l'ecrit parce que c'est ce
%  qu'on a appris a ecrire. Les vingt regles de kolonoj.py visent les
%  deux cotes a la fois — huit contre le standard (ليس، هذا، الذي،
%  سوف، ماذا، لماذا، أين، هل), cinq contre l'egyptien (ده، إزاي،
%  دلوقتي، كده، عايز), et une contre les quatre lettres پ چ ژ گ, que
%  l'arabe n'a pas : la colonne persane du releve se defendait contre
%  l'arabe, celle-ci se defend contre elle.
%
%  TROIS FORMES NE SONT DELIBEREMENT PAS RELEVEES. « كيف » est
%  levantin ET standard ; « بس » est levantin ET egyptien ; « في »
%  est l'existentiel levantin et la preposition standard. Une regle
%  sur l'une des trois crierait a chaque tableau, et un controle
%  qu'on desarme ne controle plus rien.
%
% -------------------------------------------------------------------
%  3. LA CONSTRUCTION D'ANNEXION TOMBE D'ACCORD AVEC L'IDO
% -------------------------------------------------------------------
%  La colonne gujaratie devait lire parigi.py AVANT d'ecrire une
%  ligne, parce qu'elle est a tete finale et retourne tout genitif.
%  L'ARABE EST A TETE INITIALE COMME L'IDO : « كرسي الأستاذ », la
%  chaise du maitre, pose la chaise d'abord, exactement comme
%  « stulo di la docisto ». Les couples que parigi.py signale
%  s'ecrivent donc dans l'ordre ou ils viennent, et l'outil se lit
%  APRES, pour verifier, non AVANT, pour prevoir. Deux colonnes de
%  suite, deux emplois opposes du meme outil : c'est la langue qui
%  decide a quel moment on le lit.
%
% -------------------------------------------------------------------
%  4. CE QUE CE TABLEAU-CI A TROUVE
% -------------------------------------------------------------------
%  LE POELE (46) EST « صوبيا » ET SON TUYAU (45) EST « بوري », ET LES
%  DEUX SONT TURCS. Quatre siecles d'Empire ottoman n'ont pas laisse
%  la langue de l'ecole — l'ecole est restee arabe d'un bout a
%  l'autre : أستاذ، درس، دفتر، قلم — mais ils ont meuble la piece.
%  Ce qui chauffe, ce qui fume, ce qui se pose par terre porte un nom
%  d'Istanbul.
%
%  ET L'ABAT-JOUR (tableau, alinea 15) EST « أباجورة », QUI EST
%  FRANCAIS. Le Mandat n'a dure que vingt-cinq ans contre quatre
%  siecles d'Ottomans, et il a laisse moins de mots — mais il les a
%  laisses dans la meme piece, une couche au-dessus. UNE SALLE DE
%  CLASSE LEVANTINE SE LIT COMME UNE COUPE DE TERRAIN : l'arabe au
%  fond, le turc sur le mobilier, le francais sur ce qui a ete
%  installe en dernier.
%
%  ENFIN « منيح » (19). Les bons eleves sont « التلاميذ المناح », et
%  c'est le mot que le controle protege contre l'egyptien « كويس ».
%  Le releve a maintenant trois arabes cote a cote, et le premier
%  endroit ou ils se separent est celui ou l'on dit du bien de
%  quelqu'un.
% ===================================================================

%%K t01-apar-0 apar
\VUsaut{2.0mm}
\VUcentre{12.6pt}{120}{كتيّب الشرح}
\VUsaut{3.2mm}
\VUcentre{8.4pt}{160}{تبع}
\VUsaut{2.6mm}
\VUtitre{15.6pt}{0}{لوحات دلماس المساعدة}
\VUsaut{3.4mm}
\VUornamento{9pt}
\VUsaut{5.0mm}
\VUcentre{13.2pt}{90}{السلسلة الأولى}
\VUsaut{3.0mm}
\VUfilet{16mm}
\VUsaut{5.6mm}

%%K t01-apar-1 apar
\VUtitre{15.0pt}{60}{اللوحة رقم ١}
\VUsaut{1.4mm}
\VUtitre{11.4pt}{0}{المدرسة. --- الثانوية. --- الصف.}
\VUsaut{2.2mm}
\VUfilet{20mm}
\VUsaut{6.4mm}

%%K t01-tit-1 sub
\VUpk{10.2pt}{40}{وصف عام.}
\VUsaut{3.0mm}

%%K t01-01-1 p
1. --- هاي اللوحة بتمثّل \VUgras{صف} بـ\VUgras{ثانوية}. بهالصف في تمن
\VUgras{طاولات} \textsuperscript{(18)} وتمن \VUgras{مقاعد}
\textsuperscript{(32)}. على كل مقعد بيقعدوا تلات تلاميذ.
\VUgras{الأستاذ} \textsuperscript{(7)} قاعد على \VUgras{كرسي}
\textsuperscript{(8)} بـ\VUgras{منصّته} \textsuperscript{(6)}.

%%K t01-02-1 p
2. --- عشمال المنصّة في \VUgras{خزانة} \textsuperscript{(15)}
\VUgras{إلها إزاز}. وعاليمين في \VUgras{اللوح الأسود}
\textsuperscript{(1)} مع \VUgras{المسّاحة} \textsuperscript{(2)}
و\VUgras{الطبشور} \textsuperscript{(3)}.

%%K t01-02-2 p
وفوق المنصّة معلّقين \VUgras{لوحات حيطية} \textsuperscript{(9, 11,}
\textsuperscript{12)} و\VUgras{خريطة} \textsuperscript{(10)}.

%%K t01-03-1 p
3. --- مش كل التلاميذ قاعدين بـ\VUgras{محلّهن} \textsuperscript{(17)}.
حنّا \textsuperscript{(4)} واقف عند اللوح، ماسك \VUgras{المسّاحة}
وعم يمحي \VUgras{الأشكال الهندسية}.

%%K t01-03-2 p
وبطرس جنب المنصّة، عم يلمّ \VUgras{الوظايف} \textsuperscript{(5)}
وبيجيبن للأستاذ.

%%K t01-04-1 p
4. --- \VUgras{هاد} الأستاذ هوّي أستاذ \VUgras{اللغات الحيّة}. بيعلّم
التلاميذ يحكوا ويقروا ويكتبوا لغات أجنبية \VUgras{(ألماني، إنكليزي،
إسباني، إيطالي، روسي أو فرنسي)}.

%%K t01-05-1 p
5. --- الصف كتير واسع وكتير منوّر. بآخره في \VUgras{شبّاك} عريض إلو
\VUgras{طاقتين تهوية} \textsuperscript{(78)} و\VUgras{باب إلو إزاز}،
لتهويته وتنويره.

%%K t01-05-2 p
هالصف مش بارد. بالنص في، لتدفيته، \VUgras{صوبيا} \textsuperscript{(46)}
كتير منيحة مع \VUgras{البوري} \textsuperscript{(45)} تبعها.

%%K t01-05-3 p
وعالحيط مثبّتين \VUgras{علّاقات} \textsuperscript{(58)}، عليهن معلّقين
طواقي وبالطوات التلاميذ.

%%K t01-tit-2 sub
\VUsaut{5.2mm}
\VUpk{10.2pt}{40}{ساحة اللعب}
\VUsaut{3.6mm}

%%K t01-06-1 p
6. --- \VUgras{الساعة} \textsuperscript{(63)} بتدلّ على تسعة وخمس دقايق.

%%K t01-06-2 p
\VUgras{الفرصة} \textsuperscript{(76)} هلّق خلصت والدرس عم يبلّش من جديد.
محل اللعب بـ\VUgras{الساحة} \textsuperscript{(75)}. منشوف كم تلميذ عم
يلعبوا. و\VUgras{المراقب} \textsuperscript{(74)} عم يراقبن.

%%K t01-07-1 p
7. --- وبالساحة كمان منشوف ناس غير هدول، يعني \VUgras{المفتّش}
\textsuperscript{(73)} و\VUgras{مدير الثانوية} \textsuperscript{(71)}
و\VUgras{الناظر} \textsuperscript{(72)}.

%%K t01-07-2 p
المفتّش بيفتّش المدارس، والمدير بيدير الثانوية، والناظر بيراقب الدروس.

%%K t01-07-3 p
والرابع هوّي \VUgras{الفرّاش} \textsuperscript{(77)}. حامل تحت باطه دفتر
ليسجّل الغيابات.

%%K t01-08-1 p
8. --- بآخر الساحة في \VUgras{أجهزة الرياضة} \textsuperscript{(70)}
والورشة تبع \VUgras{الأشغال اليدوية} \textsuperscript{(69)}.

%%K t01-08-2 p
وعيمين اللوحة منشوف \VUgras{صفوف} \VUgras{الموسيقى} و\VUgras{الرسم}.

%%K t01-noto-1 noto
\VUnotes{143.2mm}{%
(*) أما \textit{أسماء العمادة} فمنصح كمان بنقلها بصورتها اللاتينية.
هيك بس بتنتنّب صيغ ما إلها عدّ. وبعدين كيف بدّك تختار بين
\textit{Giovanni, Jean, Johann, Jan, Hans, John, Ivan} وغيرها ؟ منعمل
قاموس خاص لأسماء العمادة ؟ خلّينا ناخد من اللاتيني صيغة الرفع ومنقول :
\VUgras{Ioannes, Ioanna; Iulius, Iulia; Iakobus; Andreas; Lukas.}
(النحو الكامل).%
}

%%K t01-tit-3 sub
\VUpk{10.2pt}{40}{وصف مفصّل.}
\VUsaut{2.4mm}

%%K t01-tit-4 sub
\VUpk{10.2pt}{40}{التلاميذ المناح.}
\VUsaut{3.0mm}

%%K t01-09-1 p
9. --- تلاميذ المقعد الأول عيمين المنصّة هنّي \VUgras{هنري} (*)
و\VUgras{إسطفان} و\VUgras{شارل}.

%%K t01-09-2 p
هنري كتير منتبه ومجتهد، عم يسمع للأستاذ. وعلى طاولته في
\VUgras{دفتر} \textsuperscript{(28)} و\VUgras{كتاب} \textsuperscript{(29)}
و\VUgras{قاموس} \textsuperscript{(30)} و\VUgras{مقلمة}
\textsuperscript{(31)}. وكتابه فيه \VUgras{صفحات} كتير، وكل فصل فيه كم
\VUgras{فقرة}، وكل \VUgras{سطر} فيه \VUgras{كلمات} كتير.

%%K t01-09-4 p
وجاره إسطفان \textsuperscript{(24)} نشيط. عم يكتب بدفتره الدرس تبع
المرة الجاي. \VUgras{خطّه} حلو. وكتابه مسكّر، منشوف بس \VUgras{الغلاف}
\textsuperscript{(27)}. وجنبه \VUgras{الفرجار} \textsuperscript{(26)}
تبعه.

%%K t01-09-5 p
وشارل \textsuperscript{(23)} ماسك كتابه بإيدو، عم يفتحه ليعيد قراية
درسه. وإلو، متل كل تلميذ، \VUgras{دواية} \textsuperscript{(25)} قدّامه.
وهوّي مطيع.

%%K t01-10-1 p
10. --- وعالطاولة اللي جنبها في \VUgras{لويس} و\VUgras{أنطون}
و\VUgras{بولس}. لويس عم يسمّع \VUgras{درسه} \textsuperscript{(22)}
وبيجاوب على \VUgras{أسئلة} الأستاذ. وأنطون \textsuperscript{(21)} كتير
مش منتبه، وبعض المرّات الأستاذ بيعاقبه. وبولس \textsuperscript{(20)}
رفيق منيح، لطيف وبيحب يساعد. وهوّي كتير منتبه.

%%K t01-tit-5 sub
\VUsaut{4.2mm}
\VUpk{10.2pt}{40}{التلاميذ الشقيّين.}
\VUsaut{3.2mm}

%%K t01-11-1 p
11. --- ووراء هنري في \VUgras{جورج} \textsuperscript{(38)}. مش ولد
عاطل، بس \VUgras{بيحكي كتير} ومش منتبه وبيحب الشقاوة. \VUgras{طيشه}
و\VUgras{قلّة طاعته} بيعملوا \VUgras{تشويش} بالدرس وكتير مرّات بيستاهل
\VUgras{إنذار} قاسي. \VUgras{دفتر مسوّدته} \textsuperscript{(35)} وسخ،
بيبقّعه بـ\VUgras{ريشته} \textsuperscript{(34)}، ولهيك دايماً بيستعمل
\VUgras{محّايته} \textsuperscript{(36)} ؛ و\VUgras{كيس دفاتره}
\textsuperscript{(33)} مقطّع، لأنه كتير مهمل. وليش جايب
\VUgras{علبة أقلامه} \textsuperscript{(37)} عصف \VUgras{اللغات الحيّة} ؟

%%K t01-11-2 p
وجاره إدمون \textsuperscript{(39)} ما بيحبّه. صحيح هوّي شوي كسلان، بس
ساكت وهادي. ما بيحكي ؛ لكن بدل ما يسمع، بيفضّل يقرا \VUgras{القصص} اللي
بـ\VUgras{كتاب المطالعة} تبعه.

%%K t01-11-3 p
وتالت تلميذ بهالمقعد هوّي \VUgras{لودفيك} \textsuperscript{(40)}. مجتهد.
عم يكتب على \VUgras{ورقة} بيضا. وقدّامه حاطط \VUgras{ورق النشّاف}
\textsuperscript{(41)} تبعه.

%%K t01-12-1 p
12. --- تلاميذ المقعد التاني، عشمال الأستاذ، كتير زغار. الأول، الزغير
\VUgras{إميل}، لسّا ما بيعرف يكتب منيح، جايب \VUgras{لوحه}
\textsuperscript{(42)} و\VUgras{قلم لوحه} و\VUgras{إسفنجته}
\textsuperscript{(43)}.

%%K t01-12-2 p
وجنبه في \VUgras{أوغست} و\VUgras{برنار} \textsuperscript{(44)}. هاد
بيشتغل منيح، بس \VUgras{لبسه} كتير مهمل.

%%K t01-13-1 p
13. --- ووراهن في تلات \VUgras{كسالى}. \VUgras{ألبير} ما بيدرس دروسه.
و\VUgras{يعقوب} \textsuperscript{(47)} بينام بدل ما يسمع. \VUgras{كسله}
بيجيب إلو \VUgras{عتاب} وكمان \VUgras{عقوبات}. وأخيراً \VUgras{كريستيان}
\textsuperscript{(48)} كتير ما بياخد الأمور بجدّ. \VUgras{بيتصرّف} عاطل
وأبداً ما عم يحاول يصلّح حاله.

%%K t01-14-1 p
14. --- أما جيرانه فبالعكس مؤدّبين. \VUgras{إرنست} \textsuperscript{(51)}
بيحب النظام والترتيب، دايماً بيستعمل \VUgras{مسطرته}
\textsuperscript{(50)} و\VUgras{قلمه الرصاص} \textsuperscript{(49)}. وهوّي
\VUgras{خارجي}.

%%K t01-14-2 p
و\VUgras{فرنسيس} \textsuperscript{(52)} \VUgras{داخلي}، لابس البدلة
الموحّدة، بياكل وبينام بالثانوية. وهوّي \VUgras{رفيق} منيح.

%%K t01-14-3 p
وموريس عم يبري \VUgras{قلمه} بـ\VUgras{مطواته} \textsuperscript{(56)}،
وهوّي كتير مرتّب.

%%K t01-14-4 p
جايب كل كتبه، \VUgras{كتاب القواعد} \textsuperscript{(55)}
و\VUgras{دفتر ملاحظاته} \textsuperscript{(54)}، وكمان
\VUgras{المسند} \textsuperscript{(53)}. و\VUgras{شنتته}
\textsuperscript{(57)} عالأرض وراه.

%%K t01-14-5 p
والطاولتين اللي بآخر الصف قاعد عليهن \VUgras{التلاميذ المناح}
\textsuperscript{(19)}.

%%K t01-tit-6 sub
\VUsaut{4.6mm}
\VUpk{10.2pt}{40}{الرسم والموسيقى.}
\VUsaut{3.4mm}

%%K t01-15-1 p
15. --- بـ\VUgras{صف الرسم} في \VUgras{نماذج} حلوة معلّقة عالحيط
و\VUgras{لمبة} \textsuperscript{(62)} إلها \VUgras{أباجورة}، مدلّاة من
السقف. و\VUgras{الأستاذ} \textsuperscript{(60)} كتير صبور، بيصلّح
\VUgras{رسوم} \textsuperscript{(61)} التلاميذ وبيعلّمن يرسموا
\VUgras{تمثال نصفي} \textsuperscript{(59)} من الطبيعة.

%%K t01-16-1 p
16. --- و\VUgras{صف الموسيقى} بيبيّن كتير مشوّق. فيه منتعلّم نقرا
\VUgras{الموسيقى}، ونسلفج \VUgras{النوطات} \textsuperscript{(67)}
المكتوبة عالـ\VUgras{مدرّج} (دو، ري، مي، فا، صول، لا، سي\,=\,C. D. E.
F. G. A. B.)، ونعرف \VUgras{المفاتيح} ونغنّي \VUgras{ألحان} حلوة.
والنوطات بيحطّوها عالـ\VUgras{حامل} \textsuperscript{(65)}. وواحد من
التلاميذ \VUgras{عم يعزف بيانو} \textsuperscript{(64)}
و\VUgras{أستاذ الموسيقى} \textsuperscript{(66)} \VUgras{عم يوزّن}
\textsuperscript{(68)}.

%%K t01-17-1 p
17. --- ومنشوف زاوية من \VUgras{الورشة} تبع \VUgras{الأشغال اليدوية}
\textsuperscript{(69)}، وين الولاد بيقدروا يتعلّموا ينجروا الخشب
ويبردوا الحديد. هاد الشي ما بيكون بكل الثانويات.

%%K t01-tit-7 sub
\VUsaut{5.0mm}
\VUpk{10.2pt}{40}{صف العلوم.}
\VUsaut{3.6mm}

%%K t01-18-1 p
18. --- أستاذ الرياضيات بيعلّم \VUgras{التلاميذ} \VUgras{الهندسة
والحساب والعمليات} و\VUgras{النظام المتري}. ومنشوف عاللوح أهم الأشكال
الهندسية : \VUgras{دايرة} \textsuperscript{(a)}, \VUgras{مربّع}
\textsuperscript{(b)}، \VUgras{مثلّث} \textsuperscript{(c)}،
\VUgras{خط} \textsuperscript{(d)}.

%%K t01-18-2 p
ومنشوف كمان \VUgras{عملية}، مش \VUgras{جمع}، ولا \VUgras{طرح}، ولا
\VUgras{ضرب}، هيّي \VUgras{قسمة} \textsuperscript{(e)}.

%%K t01-19-1 p
19. --- وعلوحة النظام المتري، منشوف : \VUgras{المتر}
\textsuperscript{(a)}, \VUgras{الميزان} \textsuperscript{(b)}
و\VUgras{صنجاته}، و\VUgras{الليتر} \textsuperscript{(e)}،
\VUgras{الديكاليتر} \textsuperscript{(f)}, \VUgras{الديسيليتر}
و\VUgras{المتر المكعّب} \textsuperscript{(d)}.

%%K t01-19-2 p
والتلاميذ بيدرسوا كمان \VUgras{علوم الطبيعة} \textsuperscript{(11)} —
علم الحيوان \textsuperscript{(a)}، علم النبات \textsuperscript{(b)}،
علم الأرض \textsuperscript{(c)} — و\VUgras{التاريخ}
\textsuperscript{(12)}.

%%K t01-19-3 p
وبالخزانة في أجهزة \VUgras{الفيزيا} \textsuperscript{(15)}
و\VUgras{الكيميا} \textsuperscript{(16)}.

%%K t01-19-4 p
وفوق الخزانة في : \VUgras{كرة} \textsuperscript{(14)}،
\VUgras{مخروط} و\VUgras{كرة أرضية} \textsuperscript{(13)}.

%%K t01-19-5 p
وفوق اللوحات معلّقة \VUgras{خريطة} بتمثّل قارّات الدني الخمسة :
\VUgras{أوروبا} \textsuperscript{(g)}, \VUgras{آسيا}
\textsuperscript{(e)}، \VUgras{إفريقيا} \textsuperscript{(f)}،
\VUgras{أميركا} \textsuperscript{(ab)} و\VUgras{أوقيانوسيا}
\textsuperscript{(h)}، والمحيطين الكبار، \VUgras{الهادي}
\textsuperscript{(c)} و\VUgras{الأطلسي} \textsuperscript{(d)}.

\VUsaut{6.0mm}
\VUfilet{22mm}
